\subsection{The Gap Between Internal Performance and External Readiness}

Every model in this study achieved AUROC 1.00 on the internal UCI test set. By the standard metrics that dominate published clinical ML literature, all five would be described as performing excellently. When the same models were applied to MIMIC-IV, AUROC fell to values indistinguishable from chance (0.48--0.58), ECE exceeded 0.68 for every model, and conformal coverage dropped from near-target to 0.21--0.25. No model crossed a single external-validation threshold on the deployment checklist.

That gap is the central finding. It does not reflect poor algorithm design. It reflects something more basic: the UCI CKD dataset, while widely used in classification benchmarks, does not represent the distributional complexity of real clinical data. The UCI cohort was collected from a single referral hospital, carries a 62.5\% CKD prevalence, and contains clean clinical measurements that in practice are rarely all available at once. The MIMIC-IV cohort represents a broader hospital population with 23.7\% CKD prevalence and seven features completely absent from the routine EHR. Isotonic calibration that achieved ECE 0.000 internally did not survive the shift.

This pattern is consistent with findings from systematic reviews. Echouffo-Tcheugui and Kengne found that external calibration validation is the exception rather than the norm in the CKD prediction literature~\cite{EchouffoTcheugui2012}. This study provides a concrete illustration of why that gap matters: a model that looks ready is not ready.

\subsection{What Drove the External Failure}

Three factors contributed to the calibration collapse. The first is prevalence shift. A model trained and calibrated on a 62.5\% CKD population assigns probabilities near the training prevalence. The MIMIC cohort sits at 23.7\% CKD, so those probability estimates are systematically too high, producing calibration curves that fall above the diagonal at every predicted probability level. This is the most direct driver of high MIMIC ECE.

The second factor is feature missingness. Seven of the 24 features in the UCI schema are not routinely recorded in MIMIC. Urine-specific gravity, urine sugar, pus cells, pus cell clumps, bacteria, appetite, and pedal edema were all imputed using UCI training-set medians and modes. Those imputed values carry no information about individual MIMIC patients. From the model's perspective, seven features are effectively noise columns on the external cohort.

The third factor is dataset saturation. AUROC 1.00 on a 61-patient test set, with near-perfect cross-validation scores, signals that the UCI benchmark does not offer a realistic discrimination challenge. The learned decision boundaries are sharp enough to separate every patient in the training domain. Those boundaries do not generalize to a population with a different disease severity distribution and different measurement patterns.

XGB was the least poor performer externally (AUROC 0.579, ECE 0.680). That small margin appears to reflect gradient boosting's regularization limiting overfitting to UCI-specific patterns. Even so, the margin is not large enough to change the practical conclusion.

\subsection{Conformal Prediction as a Diagnostic Tool}

The conformal coverage collapse offers something calibration numbers alone do not: a direct, interpretable signal of distribution shift. The theoretical guarantee of split conformal prediction holds only when calibration and test data come from the same distribution~\cite{AngelopoulosBates2022}. Coverage of 0.22--0.25 on MIMIC, against a 0.90 target, is a quantitative statement that the MIMIC population is not exchangeable with the UCI validation set. A system operator who sees conformal coverage fall from 0.97 to 0.22 gets an immediate signal that the model is operating outside its valid domain.

That is the practical value of including conformal prediction in a deployment framework, separate from its theoretical properties. It creates a coverage audit trail. High MIMIC singleton rates for LR (1.00) and NB (0.99) might appear to suggest interpretable outputs. They do not: when coverage is only 0.24, a singleton prediction is a confident wrong answer for roughly three quarters of patients. Singleton rate without coverage context is not a useful interpretability metric.

\subsection{Comparison to Published CKD Models}

The eight externally-validated CKD models identified by Echouffo-Tcheugui and Kengne reported a range of calibration outcomes, but assessment was typically informal, often through visual calibration plots rather than ECE or Brier Score computation~\cite{EchouffoTcheugui2012}. More recent models, including the KFRE validated by Tangri and colleagues~\cite{Tangri2016} and its UK validation by Major and colleagues~\cite{Major2019}, achieve strong external discrimination (C-statistic 0.80--0.90) in prospective cohort studies with purpose-built feature ascertainment.

The AUROC values seen on MIMIC in this study (0.48--0.58) are not comparable to those results, because the KFRE and similar models were validated on cohorts where inputs were actually measured. The comparison instead reinforces a narrower point: a model applied to harmonized data with seven imputed features is not the same as a model applied to complete data. Deployment performance depends on what the deployment environment actually records, not on what the training environment provided.

\subsection{Limitations}

The MIMIC external cohort used here is the 100-patient public demonstration subset, not the full MIMIC-IV database. At 97 patients after filtering, calibration metrics carry wide uncertainty. Bootstrap confidence intervals for MIMIC ECE span ranges as large as 0.18 for some models. These results illustrate the analytic framework. They are not definitive estimates of performance in clinical deployment.

Feature harmonization introduced a form of domain leakage. The seven imputed features were filled using UCI training statistics, so the MIMIC cohort's representation of those features simply echoes the UCI training population. That is an inherent limitation of cross-dataset harmonization when features are domain-specific, and it worsens external performance in a way that cannot be fully corrected after the fact.

The UCI dataset comes from a single referral hospital in India. MIMIC-IV reflects a tertiary care academic medical center in Boston. External validation across such different contexts is useful for stress-testing but does not directly estimate performance in any specific deployment environment. Subgroup analysis was also limited: DM and HTN subgroups each fell below the 10-patient threshold, and the age subgroups (below 65: $n$\,=\,55; 65 and above: $n$\,=\,42) were small enough that bin-level ECE estimates carry meaningful uncertainty.

\subsection{Future Directions}

Three areas appear worth pursuing. First, the full MIMIC-IV database would provide several thousand patients with CKD-relevant laboratory data, allowing larger-scale calibration assessment and properly powered subgroup analysis. The extraction pipeline from this study is directly applicable to the full database.

Second, the deployment readiness framework here is intentionally simple. Criteria are binary and thresholds are set a priori without formal sample size calculations. Extending the framework to account for uncertainty in threshold exceedance, for example through bootstrap p-values for each criterion, would make the checklist scores more principled.

Third, patient-facing uncertainty display is an open design problem. Conformal prediction sets at the individual level (CKD / not-CKD / ambiguous) produce clinically interpretable labels. How clinicians and patients respond to ambiguous predictions in a real consultation, and whether explicit uncertainty communication changes treatment decisions compared to point probability outputs, has not been tested in CKD care.
